\documentclass[11pt,a4paper]{article}
\usepackage{graphicx}
\usepackage[activeacute, spanish]{babel}
\usepackage[latin1]{inputenc}
\usepackage{amssymb,amsmath}

% \let\mathnumsetfont\mathbb
\newcommand\Nset{\mathbb N}       % set of positive integer numbers
\newcommand\Zset{\mathbb Z}       % set of integer numbers
\newcommand\Rset{\mathbb R}       % set of real numbers
\newcommand\real{\mathbb R}       % set of real numbers

\author{ Guillermo T\'unez, Mat\'ias Albanesi, Mercedes Bianchi }

\title{ \Huge
        \textbf{Trabajo Pr�ctico N� 2}\\[2mm]
        \LARGE
        \textbf{Distribuci\'on el\'ectrica}\\[6mm]
        \large
        Universidad de Buenos Aires\\
        Facultad de Ciencias Exactas y Naturales\\
        Departamento de Computaci�n\\
        M�todos N�mericos\\[10mm]
        %\includegraphics[width=3cm]{uba.jpg}\\[6mm]
        Primer Cuatrimestre - 2007\\ }

\begin{document}
\maketitle
    \begin{center}
        \begin{tabular}{|c|c|c|}
            \hline
                Nombre & LU & E-mail \\
                Mat\'ias Albanesi C\'aceres & 522/00 & mac@sion.com \\
      					Mar\'ia Mercedes Bianchi & 685/97 & mbianchi9@arnet.com.ar \\      	      							
      					Guillermo T\'unez & 790/00 & guillermotunez@yahoo.com.ar\\
		
                \hline
        \end{tabular}
\end{center}


\addvspace{3cm}

\textbf{Extracto:} Este trabajo se centra en el an\'alisis de m\'etodos para resolver sistemas de ecuaciones lineales y su aplicaci\'on en problemas de flujo en grafos, en particular sobre una red de Morley.
Se desarroll\'o una variante de la eliminaci\'on gaussiana con pivoteo parcial que puede trabajar con sistemas indeterminados y se aplic\'o al problema de la red de Morley. Se estudi\'o su performance y el n\'umero de condici\'on de la matriz del sistema para estimar el error num\'erico cometido.
El algoritmo produce un resultado gr\'afico que permite comprobar f\'acilmente la correctitud de la soluci\'on y compartirla con usuarios sin que necesiten conocer detalles de la implementaci\'on.\\

\textbf{Palabras Clave:} Sistema de ecuaciones lineales, Grafo, Red de Morley, Eliminaci\'on gaussiana, Error num\'erico.

%----------------------------------------------------------------------------------------------%
\newpage

\tableofcontents

\newpage
\section{Introducci\'on}

\subsection{Sistemas de Ecuaciones Lineales}

Un sistema lineal de ecuaciones consta de n ecuaciones sobre n inc\'ognitas y puede expresarse en notaci\'on matricial como $A \times x = b$. Estas t\'ecnicas  utilizan una serie finita de operaciones aritm\'eticas para determinar la soluci\'on exacta del sistema, sujeta \'unicamente al error de redondeo.\\

Una propiedad importante acerca de la existencia o no de la soluci\'on del sistema es la existencia de una inversa $A^{-1}$ de la matriz $A$; si existe el sistema tiene soluci\'on y adem\'as es \'unica. Equivalentemente el sistema tiene soluci\'on si el determinante de la matriz $A$ no es cero. En este caso tambi\'en se dice que la matriz $A$ es no singular.\\

La soluci\'on del sistema lineal es un vector $x$ tal que $x = A^{-1} b$.

\subsection{M\'etodo de Gauss para la resoluci\'on de sistemas de ecuaciones Lineales}

El m\'etodo de Gauss, conocido tambi\'en como de triangulaci\'on o de eliminaci\'on gaussiana, nos permite resolver sistemas de ecuaciones lineales con cualquier n\'umero de ecuaciones y de inc\'ognitas.\\ 

La idea es muy simple; por ejemplo, para el caso de un sistema de tres ecuaciones con tres inc\'ognitas se trata de obtener un sistema equivalente cuya primera ecuaci\'on tenga tres inc\'ognitas, la segunda dos y la tercera una. Se obtiene as\'i un sistema triangular o en cascada.\\

Al resolver un sistema puede suprimirse, sin que var\'ie su resoluci\'on, cualquier ecuaci\'on que pueda obtenerse a partir de otras aplicando las siguientes operaciones:
\begin{itemize}
\item Producto o cociente por un n\'umero real distinto de cero.
\item Suma o diferencia de ecuaciones.
\end{itemize}

Es importante a�adir que el sistema resultante es dependiente de la forma en que apliquemos estos criterios, es decir, las ecuaciones obtenidas no son siempre las mismas, pero si lo hemos aplicado correctamente, el sistema resultante es equivalente al dado.

\newpage

\subsection{Soluci\'on de un Sistema}

La soluci\'on de un sistema es cada conjunto de valores que satisface a todas las ecuaciones.
Los tipos de sistemas pueden clasificarse en 3:\\
\begin{itemize}
\item Sistemas Compatibles
\item Sistemas Compatibles Indeterminados
\item Sistemas Incompatibles
\end{itemize}

\subsubsection{Sistemas Determinados}

Se llama sistema compatible determinado al que tiene una sola soluci\'on com\'un, es decir, que sus rectas asociadas se cortan en un s\'olo punto, por lo tanto sus vectores de direcci\'on son linealmente independientes.

\subsubsection{Sistemas Compatibles Indeterminados}

Se llama sistema compatible indeterminado al que tiene infinitas soluciones, es decir, que sus rectas asociadas son id\'enticas, por lo tanto sus ecuaciones son equivalentes.

\subsubsection{Sistemas Incompatibles}

Se llama sistema incompatible al que no tiene ninguna soluci\'on, es decir, que sus rectas asociadas son paraleleas, por lo tanto tiene alg\'un vector de direcci\'on linealmente dependiente.

\subsection{Norma Matricial}

Una norma matricial sobre el conjunto de todas las matrices de  $n \times n$ es una funci\'on de valor real, $\lVert * \rVert$ , definida en este conjunto y que satisface para todas las matrices  $A$ y $B$ de $n \times n$ y todos los n\'umeros reales $\alpha$:
\\
\\
\begin{itemize}
\item $\textbf{(i)} \lVert A \rVert \geq 0.$
\item $\textbf{(ii)} \lVert A \rVert = 0$, si y s\'olo si $A$ es $0$, la matriz con todas las entradas cero.
\item $\textbf{(iii)} \lVert \alpha A \rVert = \mid \alpha \mid \lVert A \rVert .$
\item $\textbf{(iv)} \lVert A + B \rVert \leq \lVert A \rVert + \lVert B \rVert.$
\item $\textbf{(v)} \lVert A  B \rVert \leq \lVert A \rVert  \lVert B \rVert.$
\end{itemize}

\newpage

\subsection{N\'umero de Condici\'on }

El n\'umero de condici\'on de una matriz no singular A relativo a la norma $\lVert * \rVert$ es \\
$$k(A) = \lVert A \rVert . \lVert A^-1 \rVert $$

Una matriz A es bien condicionada si $k(A)$ est\'a cerca a 1 y es  mal condicionada si $k(A)$ es significativamente mayor a 1. Dentro de este contexto, la condici\'on de la matriz es una medida de la seguridad relativa de que al tener un vector residual peque\~no implique que la soluci\'on obtenida sea cercana a la soluci\'on exacta.


%--------------------------------------------------------------------------------------------------------


\newpage
\section{Desarrollo}

\subsection{An\'alisis del problema}
	
		\subsubsection{Modelo matem\'atico}
		
		Como paso previo al dise\~no de una soluci\'on en computadora del problema, analicemos el problema y sus implicaciones.

		El problema a resolver consiste en analizar si una trayectoria de un autom\'ovil (puntos en el plano) puede completarse cumpliendo una restricci\'on (inecuaci\'on) que determina si el auto derrapa debido a una velocidad o giro excesivos.

		Este an\'alisis debe hacerse no s\'olo en los puntos dados sino en todos los puntos de la trayectoria; para esto debe plantearse un spline param\'etrico que interpole los puntos dados como dato para obtener una ecuaci\'on de la trayectoria que pueda utilizarse para controlar si se cumple la restricci\'on de no p\'erdida de adherencia.

		El intervalo de tiempo entre puntos es constante, y se lo denomina $\omega$.

		Se pide chequear si el auto derrapa en alg\'un punto mediante un m\'etodo de b\'usqueda de m\'aximos o m\'inimos de una variable.

		Por \'ultimo debe aplicarse el programa para analizar la mejor manera de recorrer un tramo de pista, o sea encontrar una serie de puntos dentro de la pista y un valor para $\omega$ tal que el auto no derrape, cuyos datos pueden encontrarse en el Enunciado.

\newpage

\subsection{An\'alisis previo al desarrollo en computadora}

		\subsubsection{Plataforma y compilador}

		La primer decisi\'on a tomar consiste en plataforma y compilador a utilizar. Como el desarrollo debe cumplir con el requisito de poder compilarse y ejecutarse con el software disponible en los laboratorios del Departamento de Computaci\'on, optamos por computadoras con procesadores compatibles con la arquitectura Intel x86, y el compilador g++ de la Free Software Foundation (disponible para Windows y Linux).

		\subsubsection{Entrada y salida}

		El formato es simple: un archivo de texto donde se especifica con una palabra clave el valor se define (por ejemplo, \textbf{omega 0.75}), uno por l\'inea, permitiendo comentarios estilo C++.
		Todos los valores son en punto flotante exceptuando \textbf{cantidad\_muestras} que es entero.
		Un ejemplo de archivo de entrada es el siguiente:

\begin{center}
	\begin{tabular}{|c|}
	\hline
Init\\
\\
masa 512\\
coeficiente\_carga 3.8851\\
coeficiente\_rozamiento 0.9389\\
omega 1.1\\
\\
cantidad\_muestras 10\\
\\
direccion\_inicial 1 -0.25\\
x,y 10 3\\
x,y 11 1\\
x,y 12 -2\\
x,y 13 -5\\
x,y 14 1\\
x,y 15 3\\
x,y 17 7\\
x,y 18 9\\
x,y 19 11\\
x,y 25 13\\
\\
End\\
	\hline
	\end{tabular}
\end{center}
		

\subsubsection{Implementaci\'on del spline param\'etrico}

		Se dise\~naron funciones (sin llegar a ser una clase C++) para calcular los coeficientes de un spline en una variable, para utilizar posteriormente dos Splines, uno para la coordenada X de la trayectoria y otro para la coordenada Y.

		La primera decisi\'on a tomar fue el tipo de frontera a utilizar en el spline: frontera libre o frontera sujeta. Optamos por utilizar una frontera sujeta en el extremo inicial de la trayectoria, ya que era de suponer que el veh\'iculo tomar\'ia la curva con una velocidad conocida, o este dato tambi\'en ser\'ia \' deseable manipularlo para analizar el \'optimo (adem\'as esto se pide por enunciado).
		Sin embargo, para el extremo final de la trayectoria optamos por una frontera libre. Esto es factible ya que el sistema de ecuaciones resultante de plantear las condiciones de este Spline es diagonal dominante y por lo tanto existe soluci\'on.

		Optamos por realizar un spline que reciba una tabla de pares ($t$, $f(t)$) en lugar de pedir que los valores de $t$ sean equidistantes, por generalidad y porque no representaba un costo alto de desarrollo.

		Tambi\'en se reciben los valores de $f_x'(t_0)$ y $f_y'(t_0)$ por par\'ametro. La funci\'on devuelve una tabla donde para cada valor de la variable $t_i$ se tienen cuatro valores correspondientes a los coeficientes del polinomio c\'ubico v\'alido para el segmento $[t_i, t_{i+1}]$.
		
		Se crearon funciones que dada la tabla y un valor cualquiera de $t$ se calcula el valor del polinomio correspondiente al segmento donde est\'a contenido $t$, y tambi\'en las derivadas primera, segunda y tercera.




\newpage
\section{Resultados}

% Pautas: deben incluir los resultados de los experimentos. utilizando el formato m�s adecuado
% para su presentaci�n. Deberan especificar claramente a que experiencia corresponde cada
% resultado. No se incluiran aqui corridas de m�quina.

% TAREAS: presentar bien los datos medidos y  aclarar a que experimento corresponde cada grafico

\subsection{C\'alculo de G(x) en distintas precisiones, en un rango amplio de valores de $\Rset$}

Se grafic\'o el exponente del resultado en punto flotante del c\'alculo de $G(x)$ en la m\'aquina.
Esto permite observar y comparar la magnitud de resultados muy grandes y muy
peque\~nos dentro del mismo gr\'afico, sacrificando poder distinguir
el signo del resultado.

\begin{center}
	\includegraphics[width=0.96\textwidth]{punto2_mantisa64_global.jpg}
	\textbf{Gr\'afico 1}: Exponente para G(x), con mantisa de 64 bits
	\label{fig:G_global_64}
\end{center}

\begin{center}
	\includegraphics[width=0.96\textwidth]{punto2_mantisa56_global.jpg}
	\textbf{Gr\'afico 2}: Exponente para G(x), con mantisa de 56 bits
	\label{fig:G_global_56}
\end{center}

\begin{center}
	\includegraphics[width=0.96\textwidth]{punto2_mantisa48_global.jpg}
	\textbf{Gr\'afico 3}: Exponente para G(x), con mantisa de 48 bits
	\label{fig:G_global_48}
\end{center}

\begin{center}
	\includegraphics[width=0.96\textwidth]{punto2_mantisa40_global.jpg}
	\textbf{Gr\'afico 4}: Exponente para G(x), con mantisa de 40 bits
	\label{fig:G_global_40}
\end{center}

\begin{center}
	\includegraphics[width=0.96\textwidth]{punto2_mantisa32_global.jpg}
	\textbf{Gr\'afico 5}: Exponente para G(x), con mantisa de 32 bits
	\label{fig:G_global_32}
\end{center}

\begin{center}
	\includegraphics[width=0.96\textwidth]{punto2_mantisa24_global.jpg}
	\textbf{Gr\'afico 6}: Exponente para G(x), con mantisa de 24 bits
	\label{fig:G_global_24}
\end{center}

\begin{center}
	\includegraphics[width=0.96\textwidth]{punto2_mantisa16_global.jpg}
	\textbf{Gr\'afico 7}: Exponente para G(x), con mantisa de 16 bits
	\label{fig:G_global_16}
\end{center}

\begin{center}
	\includegraphics[width=0.96\textwidth]{punto2_mantisa8_global.jpg}
	\textbf{Gr\'afico 8}: Exponente para G(x), con mantisa de 8 bits
	\label{fig:G_global_8}
\end{center}

\subsection{}

%----------------------------------------------------------------------------------------------%
\newpage



\newpage
\section{Discusi�n}

% PAUTAS: Se incluir� aqu� un an�lisis de los resultados obtenidos en la secci�n anterior (se
% analizara su validez, coherencia, etc). Deben analizarse como m�nimo los items pedidos en el
% enunciado. No es aceptable decir que "los recultados fueron los esperados", sin hacer clara 
% referencia a la te�ria a la cual se ajustan. Adem�s se deben mencionar los resultados
% interesantes y los casos "patol�gicos" encontrados.

Las instancias de prueba que ejecutamos no son lo sufientemente representativas como para poder extraer alguna idea del comportamiento del tiempo de ejecuci\'on en funci\'on del factor de demanda.
La complejidad de la version del algoritmo que elejimos para implementar nos acoto considerablemente los tiempos de prueba de otras instancias que nos pudieran proveer de alguna hipotesis.\\
A priori nuestra idea era que cuanto mayor sea $\beta$ mayor ser\'ia el tiempo de ejecuci\'on ya que las restricciones para los productos ser\'ian menores. Entiendase por esto al hecho de que ser\'ian menos restrictivas. Los valores de los factores de demanda multiplicados por un valor mayor a 1 har\'ia que los terminos independientes sean m\'as grandes y por ende el algoritmo tendr\'ia mayor libertad para acercentar los valores de las incognitas.\\
En el caso de que $\beta$ fuera menor a 1 el algoritmo debiera tener un comportamiento an\'alogo. Todo esto asumiendo que los valores esperados de la demanda son mayores a 1, algo que nos parece totalmente l\'ogico ya que ninguna empresa pretende vender solo un pasaje o menos de esto.\\
Las soluciones obtenidas presentan muy marcadamente dos grupos de valores. Aquellos que valen $0$ y aquellos otros que alcanzan el valor m\'aximo permitido por sus restricciones.\\



\newpage
\section{Conclusiones}

% PAUTAS: esta secci�n debe contener las conclusiones generales del trabajo. Se deben 
% mencionar las relaciones de la discuci�n sobre las que se tiene certeza, junto con 
% comentarios y observaciones generales aplicables a todo el proceso. Mencionar tambi�n 
% posibles extensiones a los m�todos, experimentos que hayan quedado pendientes, etc.

% TAREAS: conclusiones generales sobre certezas. Mejoras y cosas pendientes.


\newpage
\section{Ap�ndice A}

\subsection{Enunciado}

\parskip = 11pt
\parindent = 0pt

\begin{centering}
\bf Laboratorio de M\'etodos Num\'ericos - Primer Cuatrimestre 2007 \\
\bf Trabajo Pr\'actico N\'umero 4: Luigi s\'olo admira a las Ferraris \\
\end{centering}

\vskip 1 cm
\hrule
\vskip 0.5 cm

\textbf{Introducci\'on}

Nos encontramos en el grupo de an\'alisis num\'erico de uno de los equipos de
F\'ormula 1, y estamos a pocas semanas de un gran premio que se realizar\'a
en un aut\'odromo a\'un no visitado por la categor\'\i a. Nos enfrentamos
al serio problema de decidir la mejor forma de tomar las curvas del nuevo
circuito, y para esto solamente podemos realizar simulaciones, dado que
el fin de semana de la carrera ser\'a la primera vez que nuestros autos
tomen contacto con esta pista.

Consideremos una curva del circuito, como la que muestra la Figura 1.
Queremos que nuestro auto comience en alg\'un punto del segmento A con
una cierta velocidad inicial, y que llegue a alg\'un punto del segmento B en
buenas condiciones (es decir, a una buena velocidad final, con las cuatro
ruedas sobre el pavimento y con la carrocer\'\i a en su lugar). Nuestro
objetivo principal es que el auto pueda sortear la curva sin derrapar, dado
que a la velocidad a la que circula se pierde el control muy r\'apidamente
si el auto pierde adherencia.

\vskip 11pt

\begin{centering}

Figura 1: Ejemplo de una curva y los puntos que determinan la trayectoria. \\
\end{centering}

\vskip 11pt

Para esto, trazamos una serie de puntos por los que esperamos que pase
el auto a intervalos regulares de $\omega = 0.1$ sg. Es decir, en el
instante $t=0$ sg.~el auto estar\'a en $p_1$, en el instante
$t=0.1$ sg.~estar\'a en $p_2$, etc. Como los puntos corresponden a
instantes equiespaciados en el tiempo, entonces estar\'an m\'as separados
en los tramos en los que el auto circule a mayor velocidad, y estar\'an
m\'as juntos en los tramos en los que el auto disminuya su velocidad.

Para determinar la trayectoria
completa del auto, utilizamos un spline param\'etrico que interpole estos
puntos, con lo cual tendremos una funci\'on $T:\real\to\real^2$ que nos
indicar\'a expl\'\i citamente la posici\'on del auto en funci\'on del tiempo.
Es decir, en el instante $t$, el auto se encuentra en el punto $T(t)$. Notar
que en funci\'on de este spline tambi\'en podemos calcular el vector velocidad y
el vector aceleraci\'on en cada instante de la trayectoria.

\vfil \eject

\textbf{C\'omo decidir si el auto se sale de la trayectoria}

Sea $t$ un instante del recorrido, en el cual el auto se encuentra en el
punto $T(t)\in\real^2$. Llamemos $v(t)$ y $a(t)$ al m\'odulo del
vector velocidad y al m\'odulo del vector aceleraci\'on en el instante $t$,
respectivamente. En este instante, la fuerza de frenado que ejerce el auto
sobre el piso es $F(t) = m\: a(t)$, siendo $m$ la masa del auto. El auto
no patina en el instante $t$ si
\begin{equation}
F(t) \ \le\ (mg + v(t)^2s) \: \mu_0. \label{condicion}
\end{equation}
En esta f\'ormula, $g=9.81$ m/sg${}^2$ es la aceleraci\'on gravitatoria,
$s$ es el coeficiente de carga din\'amica (de modo tal que $v(t)^2 s\mu_0$ es
la fuerza que ejerce el aire sobre la superficie del veh\'\i culo en
direcci\'on normal a la superficie de la pista, representando la acci\'on
de los alerones), y $\mu_0$ es el coeficiente de rozamiento est\'atico
(que depende del compuesto de los neum\'aticos, de la composici\'on del
pavimento y de las temperaturas del piso y de los neum\'aticos). El auto
se mantiene dentro de la trayectoria sin derrapar si la condici\'on
(\ref{condicion}) se cumple en todos los puntos de la misma.

\textbf{Enunciado}

El objetivo del trabajo pr\'actico es implementar un programa que tome los
datos del auto y de la trayectoria desde un archivo, que calcule un
spline param\'etrico que interpole los puntos de la trayectoria y que
determine si en alg\'un punto del recorrido el auto pierde adherencia. En
caso positivo, el programa deber\'a informar el punto en el cual esto
sucede, y la velocidad y direcci\'on del auto en ese momento.

El programa deber\'a tomar los datos de entrada desde un archivo de texto
cuyo formato queda a criterio del grupo. Debe notarse que los datos de
entrada incluyen, adem\'as de los puntos de la trayectoria, los datos del
auto (masa, coeficiente de rozamiento, coeficiente de carga din\'amica, etc.)
y el intervalo $\omega$ entre puntos consecutivos de la trayectoria.

Se deber\'a utilizar un algoritmo eficiente para el c\'alculo del spline
param\'etrico. Se debe mencionar en el informe el tipo de spline que se usa
para la interpolaci\'on, comentando las diferentes opciones consideradas
y las razones que llevaron a seleccionar esta opci\'on.

El c\'alculo para determinar si el auto derrapa en alg\'un punto de la
trayectoria deber\'a estar basado en un m\'etodo num\'erico a elecci\'on del
grupo para buscar m\'aximos o m\'\i nimos de funciones de una variable.
Sugerimos consultar con los docentes antes de utilizar alg\'un m\'etodo
num\'erico que no se encuentre dentro de esta categor\'\i a.

Adem\'as de estos puntos obligatorios, los invitamos a implementar los
siguientes \'\i tems optativos que pueden ser de inter\'es para el
an\'alisis:
\begin{itemize}
\item Realizar gr\'aficos de la velocidad y la aceleraci\'on en funci\'on
del tiempo. Tambi\'en puede ser interesante graficar la fuerza $F(t)$ en
funci\'on del tiempo, para ver en qu\'e momento el giro se torna peligroso.
\item Mostrar una animaci\'on en pantalla de la trayectoria del auto, junto
con su vector velocidad y la fuerza de frenado, para ver c\'omo evolucionan
a medida que el auto recorre la pista.
\end{itemize}

\vfil \eject

\textbf{Experimentos obligatorios}

Utilizando el programa implementado, analizar la mejor forma de tomar la
curva mostrada en la Figura 2, sabiendo que $\mu_0 = 0.9389$, $s = 3.8851$
kg/m y $m = 512$ kg. Buscar la velocidad inicial que permite
tomar las curvas lo m\'as r\'apidamente posible. Analizar si conviene tomar
cada curva con un radio de giro peque\~no pero a baja velocidad, o si es mejor
describir una trayectoria m\'as abierta pero a mayor velocidad. Buscar en
qu\'e momento conviene frenar, y desde qu\'e punto ya es seguro acelerar.
El grupo que logre la trayectoria m\'as r\'apida obtendr\'a una botella
de champagne y el derecho a mojar a sus compa\~neros a modo de festejo.

\begin{centering}

Figura 2: Curva para los experimentos. \\
\end{centering}

\vskip 1.2 cm
\hrule
\vskip 0.2 cm

Fecha de entrega: Lunes 25 de Junio.
\end{document}




\newpage
\section{Ap�ndice B}

\subsection{Funciones Relevantes}

\begin{scriptsize}
\begin{verbatim}

/**
 * Triangular la matriz *this modificando tambi�n b
 * No hace el despeje. Se detiene si la matriz *this representa a un sistema incompatible
 * o indeterminado (devolviendo false).
 * Devuelve true si la matriz representa un sistema determinado con igual cantidad de
 * filas no completamente iguales a 0 que columnas.
 */
bool clase::triangular(clase& b)
{

	assert(getFil() == b.getFil());
  bool rv = true;
	// Se har� pivoteo parcial. Guardar un vector de sub�ndices para cada fila
	unsigned int nrow[getFil()];
	for (uint_t i = 0; i < getFil(); i++)
		nrow[i] = i+1; // inicialmente la fila i-�sima es la verdadera fila i
	// Triangular la matriz
	// Poner en la fila i-�sima la que tiene el elemento en la columna 
	// i-�sima con valor absoluto m�s grande
	
	for (uint_t i = 1; i <= getFil(); i++) {
    // i es la fila "l�gica" que se est� procesando
		unsigned int f = nrow[i-1]; // la fila "f�sica" usada para acceder a la matriz
		if (i > getCol()) {
			// las �ltimas filas deben estar en cero
			if (!tieneFilaEnCero(f) || !b.tieneFilaEnCero(f)) {
				rv = false; // indeterminado o incompatible
				break;
			}
		}
		else {
			// Una de las primeras filas. Buscar el elemento pivote de mayor valor absoluto,
			// haciendo intercambio de filas (un intercambio simulado, se intercambian los
			// �ndices a filas en la tabla nrow[])
			uint_t filaMejorPivote = i;
			uint_t columna = i;
			while (columna <= getCol()) {
				filaMejorPivote = i;
				f = nrow[filaMejorPivote-1];
				long double mejorPivote = abs(sub(f, columna));
				for (uint_t j = i+1; j <= getFil(); j++) {
					// Ver si el elemento sub(nrow[j-1], i) es un mejor pivote
					long double candidato = abs(sub(nrow[j-1], columna)); 
					if (candidato > mejorPivote) {
						mejorPivote = candidato;
						filaMejorPivote = j;
					}
				}
				if (mejorPivote < EPSILON) {
					// Problema: tenemos la columna i-�sima en cero a partir de la fila i en adelante
					// El sistema es indeterminado.
					// Proseguir el algoritmo, tomando como elemento pivote el de la pr�xima columna
					columna++;
				}
				else
					break; // OK, seguir el algoritmo
			}
			if (columna > getCol()) {
				rv = false; // indeterminado o incompatible
				break;
			}
			
			// Intercambiar (simuladamente) las filas si corresponde
			if (i != filaMejorPivote) {
				uint_t tmp = nrow[filaMejorPivote-1];
				nrow[filaMejorPivote-1] = nrow[i-1];
				nrow[i-1] = tmp;
			}
			
			// Para cada fila a partir de i+1 restar un m�ltiplo conveniente de la fila i
			for (uint_t j = i+1; j <= getFil(); j++) {
				unsigned int fi = nrow[i-1];
				unsigned int fj = nrow[j-1];
				if (abs(sub(fj, columna)) >= EPSILON) {
					//long double mi = sub(fj, columna) / sub(fi, columna); // multiplicador
					long double mi = sub(fj, columna); // multiplicador
					long double di = sub(fi, columna); // divisor
					sub(fj, columna) = 0.0;
					for (uint_t k = columna+1; k <= getCol(); k++) {
          			long double thissubfjk = sub(fj, k) - mi * sub(fi, k) / di; 
                // E(j) - mi*E(i) -> E(j)
                if (-EPSILON < thissubfjk && thissubfjk < EPSILON)
                			thissubfjk = 0.0;
								sub(fj, k) = thissubfjk;
          }
					for (uint_t k = 1; k <= b.getCol(); k++) {
              long double bsubfjk = b.sub(fj, k) - mi * b.sub(fi, k) / di; 
              // b.E(j) - mi*b.E(i) -> b.E(j)
              if (-EPSILON < bsubfjk && bsubfjk < EPSILON)
              		bsubfjk = 0.0;
							b.sub(fj, k) = bsubfjk;
          }
				}
				else {
					sub(fj, columna) = 0.0;
				}
			}
		}
	}

	// Aplicar los intercambios de filas hechos
	for (uint_t i = 1; i <= getFil(); i++) {
		uint_t f = nrow[i-1]; // fila real en la matriz
		// Colocar la fila real en la fila i-�sima, intercambiando si corresponde
		if (i != f) {
		
			// Intercambiar las filas i-�sima y f-�sima
			for (uint_t k = 1; k <= getCol(); k++) {
				long double tmp = sub(i, k);
				sub(i, k) = sub(f, k);
				sub(f, k) = tmp;
			}
			
			// Tambi�n intercambiar en el valor b
			for (uint_t k = 1; k <= b.getCol(); k++) {
				long double tmp = b.sub(i, k);
				b.sub(i, k) = b.sub(f, k);
				b.sub(f, k) = tmp;
			}
			
			// Indicar en la tabla nrow que esas filas fueron intercambiadas
			// Buscar cu�l elemento en la tabla nrow apuntaba a la fila f�sica i y cambiarlo a f
			for (uint_t k = i + 1; k <= getFil(); k++)
				if (nrow[k-1] == i) {
					nrow[k-1] = f; //la que antes era la fila f�sica i ahora est� en la fila f�sica f
					break;
				}
			nrow[i-1] = i;
		}
	}
	return rv;
}
\end{verbatim}
\end{scriptsize}
\newpage

\begin{scriptsize}
\begin{verbatim}
/**
 * Despejar la matriz
 * Como precondici�n, la matriz debe corresponder a un sistema triangulado superiormente,
 * de M x N, con las primeras N filas distintas de cero y las restantes (si hay) en cero.
 */
 
void clase::despejar(clase& b) {

	assert(b.getFil() == getFil()); // deben tener la misma cantidad de columnas
	assert(clase::esDeterminado(b));

	// Modificar la matriz *this y tambi�n los valores de la matriz resultado b
	for (uint_t i = getFil(); i > 0; i--) {
		if (i > getCol()) {
			// La fila debe estar toda en cero, tanto para esta matriz como para la matriz b
			assert(tieneFilaEnCero(i) && b.tieneFilaEnCero(i));
		}
		else {
		
			// Los primeros i - 1 elementos deben estar en cero
			for (uint_t j = 1; j < i; j++)
				assert(sub(i, j) == 0.0); // o usar una tolerancia
			assert(sub(i, i) != 0.0); // no puede ser cero!!!
			
			// El valor debe ser 1.0
			long double pivote = sub(i, i);
			if (pivote != 1.0) {
				sub(i, i) = 1.0;
				// Multiplicar toda la fila i-�sima de b por 1/pivote
				b.multiplicarFilaPorEscalar(i, 1.0/pivote);
			}
			
			// Ahora despejar con este valor las filas 1 .. i-1
			for (uint_t k = 1; k < i; k++) {
				pivote = sub(k, i);
				
				if (abs(pivote) >= EPSILON) {
					sub(k, i) = 0.0;
					
					// Restar de la fila k de la matriz b, la fila i multiplicada por el pivote
					b.sumarMultiploDeFila(k, -pivote, i, 1);
				}
			}
		}
	}
}
\end{verbatim}
\end{scriptsize}

\begin{scriptsize}
\begin{verbatim}
void hacerDeterminado(Matriz& A, Matriz& B, AnalisisMatriz& am) {
	assert(am.esTriangularSuperiorYDespejable && am.esIndeterminado && A.getFil() >= 
	A.getCol());
	// La matriz A corresponde a un sistema indeterminado.
	// Por lo tanto tiene algunas filas en cero, correspondientes a variables 
	// que no puede despejar.
	// Vamos a convertirlo a un sistema determinado, agregando ecuaciones de la forma
	// eje(i) = 0
	// Recorrer la matriz buscando variables sin un valor definido.
	// Sabemos que la matriz tiene am.cantFilasDistintasDeCero filas que no son todas cero.
	// Agregar las nuevas ecuaciones en las filas siguientes
	unsigned int proxFila = am.cantFilasDistintasDeCero + 1;
	unsigned int i = 1;
	for (unsigned int j = 1; j <= A.getCol(); j++) {
		// Si el valor A(i,j) == 0, tenemos la variable j-�sima indeterminada
		if (abs(A.sub(i, j)) == 0.0) {
			// Agregar una ecuaci�n de la pinta x(j-esimo) = 0
			A.sub(proxFila, j) = +1.0;
			B.multiplicarFilaPorEscalar(proxFila, 0.0); // innecesario. Si fuera distinto 
			// de cero ser�a incompatible
			// Ubicar a la siguiente fila para las ecuaciones agregadas
			proxFila++;
		} 
		else {
			// El x(j-esimo) puede despejarse. Saltar a la siguiente fila y variable
			i++;
		}
	}
}
\end{verbatim}
\end{scriptsize}
\newpage

\begin{scriptsize}
\begin{verbatim}
int calcular(InfoMatriz& im, string salidaMatriz) {
    std::ostream &log = logger::getLogger();
    // Resolver el sistema de ecuaciones
    // Triangular
    bool ok = im.matriz.triangular(im.balances);
    // Ver si el sistema es incompatible
    AnalisisMatriz am = im.matriz.analizarMatriz(im.balances);
    if (!am.esTriangularSuperiorYDespejable) {
    	im.matriz.show(log);
    	return 255;
    }
    if (am.esIncompatible) {
    	im.matriz.show(log);
    	im.balances.show(log);
    	return 196;
    }
    // Ver si el sistema es indeterminado. En este caso admite m�ltiples soluciones.
    // Asignar valores para reducirlo a un sistema determinado, y luego despejar
    if (am.esIndeterminado) {
    	hacerDeterminado(im.matriz, im.balances, am); // agregar filas a la matriz 
    	// para que pueda despejarse
    	im.matriz.show(log);
		  im.balances.show(log);
      im.matriz.triangular(im.balances); // hacer que quede triangular
    	am = im.matriz.analizarMatriz(im.balances);
    	if (!am.esDeterminado) {
    		return 128;}
    }
    // Chequear que el sistema sea determinado
    assert(am.esDeterminado);
	// Exportar la matriz para poder calcular su n�mero de condici�n con Matlab
	if(salidaMatriz != "") {
        ofstream salida_matlab(salidaMatriz.c_str());
        if (!salida_matlab.is_open() || !salida_matlab.good()) {
            return 240;}
		  // sabemos que la matriz es determinada, entonces las �ltimas filas est�n
		  // en cero. No es necesario exportarlas.
        im.matriz.show_csv(salida_matlab, 1, im.matriz.getCol(), 1, im.matriz.getCol());
        salida_matlab.close();
    }
	im.matriz.despejar(im.balances);
	im.matriz.show(log);
	im.balances.show(log);
	// Volcar el resultado al grafo, asignando direcci�n y flujo a cada eje
	set<unsigned int> ejes = im.grafo.getEjes();
	set<unsigned int>::iterator itEjes;
	for (itEjes = ejes.begin(); itEjes != ejes.end(); itEjes++) {
		// Obtener el eje
		Grafo<InfoNodo, InfoEje>::Eje eje = im.grafo.getEje(*itEjes);
		// Obtener su flujo
		long double flujoEje = im.balances.sub(*itEjes, 1);
		if (flujoEje < 0) {
			// Invertir el sentido del eje
			eje.invertirSentido();
			flujoEje = -flujoEje;
		}
		// Asignar el valor del flujo
		InfoEje infoEje;
		infoEje.flujo = flujoEje;
		eje.setInfoUsuario(infoEje);
	}
  return 0;}
\end{verbatim}
\end{scriptsize}
\newpage

\begin{scriptsize}
\begin{verbatim}
InfoMatriz leerDe(std::istream& in)
{	...
	...
	...
	 // Agregar las ecuaciones de los nodos en la matriz
		set<unsigned int> ejes = rv.grafo.getEjes();
		//set<unsigned int>::iterator it;
		for (it = ejes.begin(); it != ejes.end(); it++) {
			// Para cada nodo debe cumplirse la ecuaci�n
			// 	 SUMA(flujo saliente) - SUMA(flujo entrante) = balance
			// Si para un nodo un eje es saliente, se agrega un +1 a la matriz; 
			// eso significa que el eje est� sumandole al balance del nodo.
			// En cambio, si el eje es entrante hay que agregarlo como -1 a la matriz
			Grafo<InfoNodo, InfoEje>::Eje eje = rv.grafo.getEje(*it);
			rv.matriz.sub(eje.getNodoOrigen(), *it) = +1.0; // es saliente del nodo origen
			rv.matriz.sub(eje.getNodoDestino(), *it) = -1.0; // es entrante al nodo destino
		}
  	// Agregar las restricciones de Morley en la matriz
		// Las restricciones de Morley dicen: el flujo en cada eje del backbone es igual.
		// En nuestra implementaci�n exigimos adem�s que todos los ejes del backbone tengan
		// el mismo sentido y el backbone sea lineal, para poder resolver el problema con un
		// sistema de ecuaciones planteado en una matriz.
		// Planteadas como ecuaciones, las restricciones son:
		//     eje(i) == eje(j)
		// y para llevar a una matriz:
		//     eje(i) - eje(j) = 0
		int fila = rv.grafo.getCantNodos(); // colocar a partir de esta fila las 
		// ecuaciones de la restricci�n de Morley
		if (rv.listaEjesBackbone.size() >= 2) {
			list<unsigned int>::iterator itLista = rv.listaEjesBackbone.begin();
			unsigned int prevEjeId = *itLista; // como la lista tiene al menos un 
			// elemento, esto no se pincha
			for (itLista++; itLista != rv.listaEjesBackbone.end(); itLista++) {
				unsigned int proxEjeId = *itLista;
				fila++;
				rv.matriz.sub(fila, prevEjeId) = +1.0;
				rv.matriz.sub(fila, proxEjeId) = -1.0;
				rv.balances.sub(fila, 1) = 0.0; // innecesario (balances se inicializa a 0)
				// preparar para la siguiente ecuaci�n
				prevEjeId = proxEjeId;
			}
}
\end{verbatim}
\end{scriptsize}



\newpage
\section{Ap�ndice C}

\subsection{Formato archivo de entrada}

Explicar formato del archivo de entrada!!!

Enunciado:

El programa deber\'a tomar los datos de entrada desde un archivo de texto
cuyo formato queda a criterio del grupo. Debe notarse que los datos de
entrada incluyen, adem\'as de los puntos de la trayectoria, los datos del
auto (masa, coeficiente de rozamiento, coeficiente de carga din\'amica, etc.)
y el intervalo $\omega$ entre puntos consecutivos de la trayectoria.


\end{document}

